\subsection{Structural Principles and Projective Regimes}
  \label{subsec:structural-principles-and-projective-regimes}

  All effective physical observables arise through a relational
  projection~$\Pi$ from the underlying configuration space~$\Omega$ of
  the~$\chi$ substrate.
  Effective descriptions are constrained by three structural principles that
  govern the admissibility of projected regimes.

  \paragraph{Principle I: Substratic Locality and Bounded Relaxation.}
    The admissibility constraints encoded in~$\chi$ are strictly local and
    subject to universal bounds.
    The transport of relational structure within the effective reprojection
    sequence admits a maximal admissible flux, constraining the rate at which
    structural information can be redistributed.
    This bound arises from the intrinsic stability conditions of~$\chi$ itself.

  \paragraph{Principle II: Non-Injective Projective Realization.}
    The projection $\Pi : \Omega \rightarrow O$ from substratic configurations
    to effective observables is generically non-injective.
    Distinct configurations of~$\chi$ may be structurally identified at the
    level of observable descriptions.
    This implies that effective descriptions need not admit a factorisable or
    locally complete representation, even when the underlying dynamics remains
    strictly local.

    \begin{remark}[Non-injectivity as structural necessity]
Principle~II is not a modelling hypothesis but a provable structural necessity.
Any surjective projection $\Pi : \Omega \to O$ whose observables are entirely
defined by $\Pi$ and whose effective level is not isomorphic to $\Omega$ must
be non-injective, with $S_\Pi > 0$.
A fully injective emergent description would reduce emergence to a relabelling
of the fundamental level --- a contradiction.
Within this projective-emergent setting, the following equivalences hold:
\begin{equation}
  \text{genuine emergence}
  \;\Longleftrightarrow\;
  \Pi \text{ non-injective}
  \;\Longleftrightarrow\;
  S_\Pi > 0.
\end{equation}
A complete proof, independent of any specific physical framework, is given
in~\cite{Beau2026l}; the structural consequences for gauge structure and
Bell correlations are developed in Appendix~B.23 and~\cite{Beau2026b}.
    \end{remark}

  \paragraph{Principle III: Projective Compensation.}
    Whenever $\Pi$ fails to resolve the full relational complexity of~$\chi$,
    effective descriptions compensate through the inflation of effective
    parameters---temperature, curvature, horizon structure---encoding
    unresolved relational structure within a reduced descriptive framework.
    These quantities function as Lagrange multipliers, not as additional
    fundamental degrees of freedom.

    Together, these principles delineate distinct projective regimes.
    When the projection is approximately injective and the effective
    reprojection sequence is far from saturation, standard local and geometric
    descriptions apply.
    Near structural saturation or strong non-injectivity, effective parameters
    grow large, signaling the breakdown of spacetime-based descriptions.
